\documentclass[10pt,a4paper]{amsart}
\usepackage{amsmath,amssymb,amsthm,graphicx}
% Retained preamble for the Stacks Project chapter {CH} (shared project preamble).
% The chapter file pinned-original-source.tex only carries \input{{preamble}} on its first line,
% so no preamble text lives inside the chapter itself.  The Stacks Project uses one single shared
% preamble for every chapter: the snapshot stages it three times (sources/stacks-chow/preamble.tex,
% sources/stacks-etale-cohomology/preamble.tex and focused-seeds/AGP/sources/stacks-snapshot-20261011/
% preamble.tex) and those copies are identical in size.  The block below is copied mechanically,
% line for line, from sources/stacks-chow/preamble.tex lines 14-18 and 154-217: the Xy-pic setup and
% the native theorem declarations together with the mathematical macros.  Nothing was added, removed
% or redefined; the wrapper class and the amsmath/amssymb/amsthm/graphicx base packages are supplied
% by the assembler.
% For commutative diagrams we use Xy-pic
% Chapter: stacks-cohomology.tex -- Cohomology of Sheaves.
\usepackage[all]{xy}

% We use 2cell for 2-commutative diagrams.
\xyoption{2cell}

% Theorem environments.
%
\theoremstyle{plain}
\newtheorem{theorem}[subsection]{Theorem}
\newtheorem{proposition}[subsection]{Proposition}
\newtheorem{lemma}[subsection]{Lemma}

\theoremstyle{definition}
\newtheorem{definition}[subsection]{Definition}
\newtheorem{example}[subsection]{Example}
\newtheorem{exercise}[subsection]{Exercise}
\newtheorem{situation}[subsection]{Situation}

\theoremstyle{remark}
\newtheorem{remark}[subsection]{Remark}
\newtheorem{remarks}[subsection]{Remarks}

\numberwithin{equation}{subsection}

% Macros
%
\def\lim{\mathop{\mathrm{lim}}\nolimits}
\def\colim{\mathop{\mathrm{colim}}\nolimits}
\def\Spec{\mathop{\mathrm{Spec}}}
\def\Hom{\mathop{\mathrm{Hom}}\nolimits}
\def\Ext{\mathop{\mathrm{Ext}}\nolimits}
\def\SheafHom{\mathop{\mathcal{H}\!\mathit{om}}\nolimits}
\def\SheafExt{\mathop{\mathcal{E}\!\mathit{xt}}\nolimits}
\def\Sch{\mathit{Sch}}
\def\Mor{\mathop{\mathrm{Mor}}\nolimits}
\def\Ob{\mathop{\mathrm{Ob}}\nolimits}
\def\Sh{\mathop{\mathit{Sh}}\nolimits}
\def\NL{\mathop{N\!L}\nolimits}
\def\CH{\mathop{\mathrm{CH}}\nolimits}
\def\proetale{{pro\text{-}\acute{e}tale}}
\def\etale{{\acute{e}tale}}
\def\QCoh{\mathit{QCoh}}
\def\Ker{\mathop{\mathrm{Ker}}}
\def\Im{\mathop{\mathrm{Im}}}
\def\Coker{\mathop{\mathrm{Coker}}}
\def\Coim{\mathop{\mathrm{Coim}}}

% Boxtimes
%
\DeclareMathSymbol{\boxtimes}{\mathbin}{AMSa}{"02}

%
% Macros for moduli stacks/spaces
%
\def\QCohstack{\mathcal{QC}\!\mathit{oh}}
\def\Cohstack{\mathcal{C}\!\mathit{oh}}
\def\Spacesstack{\mathcal{S}\!\mathit{paces}}
\def\Quotfunctor{\mathrm{Quot}}
\def\Hilbfunctor{\mathrm{Hilb}}
\def\Curvesstack{\mathcal{C}\!\mathit{urves}}
\def\Polarizedstack{\mathcal{P}\!\mathit{olarized}}
\def\Complexesstack{\mathcal{C}\!\mathit{omplexes}}
% \Pic is the operator that assigns to X its picard group, usage \Pic(X)
% \Picardstack_{X/B} denotes the Picard stack of X over B
% \Picardfunctor_{X/B} denotes the Picard functor of X over B
\def\Pic{\mathop{\mathrm{Pic}}\nolimits}
\def\Picardstack{\mathcal{P}\!\mathit{ic}}
\def\Picardfunctor{\mathrm{Pic}}
\def\Deformationcategory{\mathcal{D}\!\mathit{ef}}

% --- appended 2026-10-12 (Chief): the project own definitions of the reference/slogan/history
% display environments, copied mechanically from the shared Stacks preamble lines 8-12
% (\usepackage{verbatim} plus the three \newenvironment lines, which the chapter preamble keeps
% verbatim from the project preamble). Nothing added, removed or redefined: these environments
% carry no mathematical content and the project itself typesets them as comments. Without these
% three lines a counted statement whose own text uses them fails with Environment slogan undefined.
\usepackage{verbatim}
\newenvironment{reference}{\comment}{\endcomment}
\newenvironment{slogan}{\comment}{\endcomment}
\newenvironment{history}{\comment}{\endcomment}
\usepackage{cleveref}
\crefname{originallemma}{original lemma}{original lemmas}
\Crefname{originallemma}{Original lemma}{Original lemmas}
\makeatletter
\newlabel{lemma-cech-Hausdorff-quasi-compact}{{09V2}{source}{}{sourcepointer-1}{}}
\newlabel{lemma-cech-Hausdorff-quasi-compact@cref}{{[originallemma][0][]09V2}{[1][0][]source}{}{}{}}
\makeatother
\crefname{originallemma}{original lemma}{original lemmas}
\Crefname{originallemma}{Original lemma}{Original lemmas}
\makeatletter
\newlabel{lemma-injective-flasque}{{09SX}{source}{}{sourcepointer-2}{}}
\newlabel{lemma-injective-flasque@cref}{{[originallemma][0][]09SX}{[1][0][]source}{}{}{}}
\makeatother
\crefname{originallemma}{original lemma}{original lemmas}
\Crefname{originallemma}{Original lemma}{Original lemmas}
\makeatletter
\newlabel{lemma-injective-trivial-cech}{{01EP}{source}{}{sourcepointer-3}{}}
\newlabel{lemma-injective-trivial-cech@cref}{{[originallemma][0][]01EP}{[1][0][]source}{}{}{}}
\makeatother
\crefname{originallemma}{original lemma}{original lemmas}
\Crefname{originallemma}{Original lemma}{Original lemmas}
\makeatletter
\newlabel{topology-lemma-lift-covering-of-quasi-compact-hausdorff-subset}{{0CY5}{source}{}{sourcepointer-4}{}}
\newlabel{topology-lemma-lift-covering-of-quasi-compact-hausdorff-subset@cref}{{[originallemma][0][]0CY5}{[1][0][]source}{}{}{}}
\makeatother
\crefname{originallemma}{original lemma}{original lemmas}
\Crefname{originallemma}{Original lemma}{Original lemmas}
\makeatletter
\newlabel{topology-lemma-refine-covering}{{09UW}{source}{}{sourcepointer-5}{}}
\newlabel{topology-lemma-refine-covering@cref}{{[originallemma][0][]09UW}{[1][0][]source}{}{}{}}
\makeatother
\title{Cohomology of a quasi-compact Hausdorff subset as the colimit of the cohomologies of its open neighbourhoods}
\author{The Stacks Project Authors}
\date{stacks-project:cohomology.tex:local-snapshot-20261011; GNU FDL 1.2 or later}
\begin{document}
\maketitle
\noindent\textit{Editorial note.} All delivered mathematical text is the author's own, copied from the pinned Stacks Project chapter file cohomology.tex; no statement or proof step was rewritten, reordered or removed.\par\medskip
\noindent\textit{Editorial note.} This card was re-admitted after a gate audit: it had been excluded only because the chapter's retained preamble previously lacked the reference/slogan/history environments. The shared preamble now carries them (appended 2026-10-12), so the card uses that file as its preamble with no per-card copy.\par\medskip
\noindent\textit{Editorial note.} The proof contains the author's sentence 'This follows using short exact sequences and dimension shifting; details omitted.' at the reduction to injectives; the sentence is reproduced verbatim, and the rest of the argument is written out in full. The author's parenthetical offering an alternative to the appeal to flasqueness is likewise reproduced.\par\medskip
\noindent\textit{Editorial note.} The author's own bibliographic remark in a reference environment at the head of the statement is preserved in the source span and is not typeset, because the project's own preamble defines reference environments as comments; no bibliography excerpt is needed.\par\medskip
\noindent\textit{Editorial note.} The label lemma-cohomology-of-closed is defined only in cohomology.tex, so no cross-chapter restatement had to be weighed for this card.\par\medskip
\noindent\textit{Allowed standard prerequisites (editorial).}\begin{itemize}
\item Lifting a covering of a quasi-compact Hausdorff subset to a refinement by opens of the ambient space. Topology, Lemma topology-lemma-lift-covering-of-quasi-compact-hausdorff-subset, invoked at proof lines 2874-2877, 2922-2925 and 2935-2938. Supplies the shrinkage steps that make the lifts glue and the lifted cochain a cocycle.
\item Refinement of an open covering of a quasi-compact space. Topology, Lemma topology-lemma-refine-covering, invoked at proof lines 2909-2911. Supplies the refinement after which every Cech cochain value lies in the presheaf of extendable sections.
\item Cech cohomology compares with cohomology for a quasi-compact Hausdorff subset. Same chapter, Lemma lemma-cech-Hausdorff-quasi-compact, invoked at proof lines 2888-2891. Supplies that a cohomology class comes from a Cech class of a covering of Z.
\item Injectives are flasque and have trivial higher Cech cohomology. Same chapter, Lemmas lemma-injective-flasque and lemma-injective-trivial-cech, invoked at proof lines 2916 and 2943-2945. Supplies the lifting of sections and the triviality of the class inside Cech cohomology.
\end{itemize}\medskip
\ifdefined\hypertarget\hypertarget{sourcepointer-1}{}\fi
\noindent\textit{Source reference.} Original lemma 09V2: sources/stacks-cohomology/pinned-original-source.tex, source lines 2667--2667. Pointer outside the counted span, invoked once for context; the load-bearing argument of the counted statement is the author's proof reproduced in full in the proof span.\par\medskip
\ifdefined\hypertarget\hypertarget{sourcepointer-2}{}\fi
\noindent\textit{Source reference.} Original lemma 09SX: sources/stacks-cohomology/pinned-original-source.tex, source lines 1903--1903. Pointer outside the counted span, invoked once for context; the load-bearing argument of the counted statement is the author's proof reproduced in full in the proof span.\par\medskip
\ifdefined\hypertarget\hypertarget{sourcepointer-3}{}\fi
\noindent\textit{Source reference.} Original lemma 01EP: sources/stacks-cohomology/pinned-original-source.tex, source lines 1408--1408. Pointer outside the counted span, invoked once for context; the load-bearing argument of the counted statement is the author's proof reproduced in full in the proof span.\par\medskip
\ifdefined\hypertarget\hypertarget{sourcepointer-4}{}\fi
\noindent\textit{Source reference.} Original lemma 0CY5: sources/stacks-chow/tags.txt, source lines 16230--16230. Pointer outside the counted span, invoked once for context; the load-bearing argument of the counted statement is the author's proof reproduced in full in the proof span.\par\medskip
\ifdefined\hypertarget\hypertarget{sourcepointer-5}{}\fi
\noindent\textit{Source reference.} Original lemma 09UW: sources/stacks-chow/tags.txt, source lines 12439--12439. Pointer outside the counted span, invoked once for context; the load-bearing argument of the counted statement is the author's proof reproduced in full in the proof span.\par\medskip
\begin{lemma}
\label{lemma-cohomology-of-closed}
\begin{reference}
\cite[Expose V bis, 4.1.3]{SGA4}
\end{reference}
Let $X$ be a topological space. Let $Z \subset X$ be a quasi-compact subset
such that any two points of $Z$ have disjoint open neighbourhoods in $X$.
For every abelian sheaf $\mathcal{F}$ on $X$ the canonical
map
$$
\colim H^p(U, \mathcal{F})
\longrightarrow
H^p(Z, \mathcal{F}|_Z)
$$
where the colimit is over open neighbourhoods $U$ of $Z$ in $X$
is an isomorphism.
\end{lemma}

\begin{proof}
We first prove this for $p = 0$. Injectivity follows from
the definition of $\mathcal{F}|_Z$ and holds in general
(for any subset of any topological space $X$). Next, suppose that
$s \in H^0(Z, \mathcal{F}|_Z)$. Then we can find opens $U_i \subset X$
such that $Z \subset \bigcup U_i$ and such that $s|_{Z \cap U_i}$
comes from $s_i \in \mathcal{F}(U_i)$. It follows that
there exist opens $W_{ij} \subset U_i \cap U_j$ with
$W_{ij} \cap Z = U_i \cap U_j \cap Z$ such that
$s_i|_{W_{ij}} = s_j|_{W_{ij}}$. Applying
Topology, Lemma
\ref{topology-lemma-lift-covering-of-quasi-compact-hausdorff-subset}
we find opens $V_i$ of $X$ such that $V_i \subset U_i$ and
such that $V_i \cap V_j \subset W_{ij}$. Hence we see that
$s_i|_{V_i}$ glue to a section of $\mathcal{F}$ over the
open neighbourhood $\bigcup V_i$ of $Z$.

\medskip\noindent
To finish the proof, it suffices to show that if $\mathcal{I}$ is an
injective abelian sheaf on $X$, then $H^p(Z, \mathcal{I}|_Z) = 0$
for $p > 0$. This follows using short exact sequences and dimension
shifting; details omitted. Thus, suppose $\overline{\xi}$ is an element
of $H^p(Z, \mathcal{I}|_Z)$ for some $p > 0$.
By Lemma \ref{lemma-cech-Hausdorff-quasi-compact}
the element $\overline{\xi}$ comes from
$\check{H}^p(\mathcal{V}, \mathcal{I}|_Z)$
for some open covering $\mathcal{V} : Z = \bigcup V_i$ of $Z$.
Say $\overline{\xi}$ is the image of the class of a cocycle
$\xi = (\xi_{i_0 \ldots i_p})$ in
$\check{\mathcal{C}}^p(\mathcal{V}, \mathcal{I}|_Z)$.

\medskip\noindent
Let $\mathcal{I}' \subset \mathcal{I}|_Z$ be the subpresheaf
defined by the rule
$$
\mathcal{I}'(V) =
\{s \in \mathcal{I}|_Z(V) \mid
\exists (U, t),\ U \subset X\text{ open},
\ t \in \mathcal{I}(U),\ V = Z \cap U,\ s = t|_{Z \cap U} \}
$$
Then $\mathcal{I}|_Z$ is the sheafification of $\mathcal{I}'$.
Thus for every $(p + 1)$-tuple $i_0 \ldots i_p$ we can find an
open covering $V_{i_0 \ldots i_p} = \bigcup W_{i_0 \ldots i_p, k}$
such that $\xi_{i_0 \ldots i_p}|_{W_{i_0 \ldots i_p, k}}$ is
a section of $\mathcal{I}'$. Applying
Topology, Lemma \ref{topology-lemma-refine-covering}
we may after refining $\mathcal{V}$ assume that each
$\xi_{i_0 \ldots i_p}$ is a section of the presheaf $\mathcal{I}'$.

\medskip\noindent
Write $V_i = Z \cap U_i$ for some opens $U_i \subset X$.
Since $\mathcal{I}$ is flasque (Lemma \ref{lemma-injective-flasque})
and since $\xi_{i_0 \ldots i_p}$ is a section of $\mathcal{I}'$
for every $(p + 1)$-tuple $i_0 \ldots i_p$ we can choose
a section $s_{i_0 \ldots i_p} \in \mathcal{I}(U_{i_0 \ldots i_p})$
which restricts to $\xi_{i_0 \ldots i_p}$ on
$V_{i_0 \ldots i_p} = Z \cap U_{i_0 \ldots i_p}$.
(This appeal to injectives being flasque can be avoided by an
additional application of
Topology, Lemma
\ref{topology-lemma-lift-covering-of-quasi-compact-hausdorff-subset}.)
Let $s = (s_{i_0 \ldots i_p})$ be the corresponding cochain
for the open covering $U = \bigcup U_i$.
Since $\text{d}(\xi) = 0$ we see that the sections
$\text{d}(s)_{i_0 \ldots i_{p + 1}}$ restrict to zero
on $Z \cap U_{i_0 \ldots i_{p + 1}}$. Hence, by the initial
remarks of the proof, there exists open subsets
$W_{i_0 \ldots i_{p + 1}} \subset U_{i_0 \ldots i_{p + 1}}$
with $Z \cap W_{i_0 \ldots i_{p + 1}} = Z \cap U_{i_0 \ldots i_{p + 1}}$
such that $\text{d}(s)_{i_0 \ldots i_{p + 1}}|_{W_{i_0 \ldots i_{p + 1}}} = 0$.
By Topology, Lemma
\ref{topology-lemma-lift-covering-of-quasi-compact-hausdorff-subset}
we can find $U'_i \subset U_i$ such that $Z \subset \bigcup U'_i$
and such that $U'_{i_0 \ldots i_{p + 1}} \subset W_{i_0 \ldots i_{p + 1}}$.
Then $s' = (s'_{i_0 \ldots i_p})$ with
$s'_{i_0 \ldots i_p} = s_{i_0 \ldots i_p}|_{U'_{i_0 \ldots i_p}}$
is a cocycle for $\mathcal{I}$ for the open covering
$U' = \bigcup U'_i$ of an open neighbourhood of $Z$.
Since $\mathcal{I}$ has trivial higher {\v C}ech cohomology groups
(Lemma \ref{lemma-injective-trivial-cech})
we conclude that $s'$ is a coboundary. It follows that the image of
$\xi$ in the {\v C}ech complex for the open covering
$Z = \bigcup Z \cap U'_i$ is a coboundary and we are done.
\end{proof}

\end{document}
